% B05/A023: retained-U payment for fully overlapping bent one-point surgery.
\section*{A retained-$U$ payment for fully overlapping bent surgery}

All logarithms are base two.  For a Boolean function $f$ write
\[
 U(f)=\sum_S\widehat f(S)^2
 \log\frac{\sum_{T\supseteq S}\widehat f(T)^2}{\widehat f(S)^2},
\]
with the usual zero-mass convention.

\subsection*{Exact shared-variable switch identity}

Let $g,h:\{-1,1\}^n\to\{-1,1\}$ be arbitrary Boolean functions on the
same physical variables.  Set
\[
 a_T=\widehat g(T),\qquad b_T=\widehat h(T),\qquad d_T=b_T-a_T,
\]
and define their upper masses
\[
 A_T=\sum_{R\supseteq T}a_R^2,\quad
 B_T=\sum_{R\supseteq T}b_R^2,\quad
 D_T=\sum_{R\supseteq T}d_R^2.
\]
Let $z$ be a fresh $d$-bit selector block, let $p=2^{-d}$, and put
\[
 F(z,x)=\begin{cases}h(x),&z=z_\star,\\g(x),&z\ne z_\star.\end{cases}       \tag{1}
\]
Then
\begin{align}
 U(F)={}&\sum_T c_T^2\log\frac{(1-p)A_T+pB_T}{c_T^2}
 +p(1-p)\sum_Td_T^2\log\frac{D_T}{d_T^2} \\
 &+\left(\frac{pd}{2}-p^2d\right)\sum_Td_T^2,                 \tag{2}
\end{align}
where $c_T=(1-p)a_T+pb_T$.

Indeed, the control-empty coefficient is $c_T$, while every nonempty
control coefficient has magnitude $p|d_T|$.  At control frequency $Q$ the
upper mass is
\[
 \begin{cases}
 (1-p)A_T+pB_T,&Q=\varnothing,\\
 p2^{-|Q|}D_T,&Q\ne\varnothing.
 \end{cases}
\]
Summing $d-|Q|$ over nonempty $Q$ gives $d/(2p)-d$, proving (2).  No
orthogonality or disjointness is used.

The complete retained-child charge is
\[
 J=U(F)-(1-p)U(g)-pU(h).                                  \tag{3}
\]
If $\delta=\Pr[g\ne h]$, then each selector coordinate is physically
pivotal with probability $2p\delta$, so
\[
 I_{\rm ctl}(F)=2pd\delta.                                \tag{4}
\]

\subsection*{Bent one-point surgery}

\paragraph{Theorem.}
Let $k\ge8$ be even, let $g$ be any bent Boolean function on $k$ bits, and
let $h$ be obtained from $g$ by flipping its value at one input point.
For every $d\ge1$, the shared-variable switch (1) satisfies
\[
 \boxed{J\le71 I_{\rm ctl}(F).}                           \tag{5}
\]
The bound is independent of $k$, $d$, the bent function, and the flipped
point.  The two children share every physical data variable and have
completely overlapping Fourier supports.

\paragraph{Proof.}
Put $N=2^k$, $\alpha=N^{-1/2}$ and $\delta=N^{-1}=\alpha^2$.  Bentness gives
$a_T=\alpha s_T$ with $s_T\in\{-1,1\}$.  A one-point flip gives
\[
 b_T-a_T=-2\delta t_T,qquad t_T\in\{-1,1\}.              \tag{6}
\]
In particular $\Pr[g\ne h]=\delta$.

For each $T$, let $L_T=k-|T|$ and introduce the dimensionless quantities
\[
 \frac{B_T}{A_T}=1+u_T,qquad
 \frac{b_T^2}{\delta}=1+v_T,qquad
 \frac{c_T^2}{\delta}=1+w_T.                             \tag{7}
\]
Since $A_T=2^{-|T|}$ and the upper-cone average of the signs $s_Rt_R$ lies
in $[-1,1]$,
\[
 |u_T|,|v_T|\le\frac{17}{4}\alpha,qquad
 |w_T|\le\frac{33}{8}p\alpha                            \tag{8}
\]
for $k\ge8$.  Moreover
\[
 v_T=-4\alpha s_Tt_T+4\delta,qquad
 w_T=-4p\alpha s_Tt_T+4p^2\delta,                        \tag{9}
\]
so
\[
 pv_T-w_T=4p(1-p)\delta.                                 \tag{10}
\]

The $L_T$ terms in (2)--(3) cancel exactly, atom by atom.  After converting
temporarily to natural logarithms, the remaining scalar term is
\[
 \Phi_p(u,v,w)
 =(1+w)\log\frac{1+pu}{1+w}
 -p(1+v)\log\frac{1+u}{1+v}.                             \tag{11}
\]
Thus
\[
 J=\frac{\delta}{\ln2}\sum_T\Phi_p(u_T,v_T,w_T)
 +(1-2p)I_{\rm ctl}(F).                                  \tag{12}
\]

At the origin, $\Phi_p=0$ and
\[
 \nabla\Phi_p(0,0,0)=(0,p,-1).
\]
On the box (8), using $p\le1/2$ and $\alpha\le1/16$, direct
differentiation gives
\[
 |\Phi_{uu}|\le4p,quad |\Phi_{uv}|\le\tfrac32p,quad
 |\Phi_{uw}|\le\tfrac65p,quad |\Phi_{vv}|\le\tfrac32p,quad
 |\Phi_{ww}|\le\tfrac65,quad \Phi_{vw}=0.              \tag{13}
\]
Taylor's theorem with integral remainder, (8), and (10) therefore give
\begin{align*}
 \Phi_p(u_T,v_T,w_T)
 &\le4p\delta+\frac12p\delta\left[
 4\left(\frac{17}{4}\right)^2
 +3\left(\frac{17}{4}\right)^2
 +\frac65\frac{17}{4}\frac{33}{8}
 +\frac32\left(\frac{17}{4}\right)^2
 +\frac35\left(\frac{33}{8}\right)^2\right]\\
 &<97p\delta.                                             \tag{14}
\end{align*}
There are $N=1/\delta$ Fourier atoms.  Hence (12), $97/\ln2<140$, and
$I_{\rm ctl}=2pd\delta$ imply
\[
 J\le140p\delta+I_{\rm ctl}
 \le\left(\frac{70}{d}+1\right)I_{\rm ctl}
 \le71I_{\rm ctl},
\]
which proves (5).\qed

\paragraph{Payment scope.}
The selector stopping rule is physical, adaptive, and finite: reveal the
$d$ fresh controls until an inactive bit appears, otherwise reach the
flipped child.  Equation (3) retains the two child packets with their exact
physical probabilities, and (5) spends each selector influence once.  The
theorem removes disjointness and Fourier orthogonality for an unbounded
fully overlapping family.  It does not cover repeated coherent surgeries
through an arbitrary hierarchy; controlling the accumulation of several
changing overlap patterns remains open.
